\section{Introduction}
\label{sec:introduction}

Most computational histories are collected to answer questions already known
at capture time: which call failed, which state changed, which action produced
a reward, or which artifact was released. The schema names the objects in
advance. Later work searches, aggregates, or predicts within that object world.

Real work creates a harder failure mode. A team can preserve every field it
knows and still lose the relation that a future participant needs, because the
relation did not yet have a name. Permission can be mistaken for accepted
responsibility. An executed action can be mistaken for a settled obligation.
A collection of verified components can be mistaken for the one project state
a successor is authorized to continue. The missing information is not always a
missing value. It may be a missing object boundary.

This paper asks whether historical infrastructure can preserve a capacity to
repair that kind of omission. The proposal builds on three Kung Fu Decisions:
KFD-4 binds views to declared perspectives and permits source-preserving replay
and transformation; KFD-5 separates candidate genesis from fact-bound
qualification; draft KFD-6 proposes a bounded loop that compares generation
methods over causal experience without allowing a generator to certify or
promote itself \cite{kfd4Decision2026,kfd5Decision2026,kfd6Decision2026}.

The central thesis is conditional:

\begin{quote}
A well-formed causal Episode corpus can preserve more than prior answers. It
can preserve a governed option to replay decisions, expose unnamed burdens,
qualify candidate Primitives, and revise the ontology under which future
questions are formed.
\end{quote}

The word \emph{option} matters. A corpus does not discover merely by existing.
A generated label is not a Primitive. A plausible retrospective story is not
causal evidence. The proposed value appears only when later participants can
identify an immutable evidence cut, reconstruct perspectives with visible
loss, compare plural methods and conservative baselines, reject attractive
false candidates, and keep promotion authority separate from generation.

The contribution is distinct from two adjacent KFD papers. The Foundation
paper develops facts, trust, cooperation, and boundary primitives as an ordered
systems model \cite{kfdFoundationPaper2026}. The Observer paper develops
declared timelines and perspective-preserving replay
\cite{kfdObserverPaper2026}. This paper treats replayable Episodes as a
long-lived input to ontology revision and focuses on the KFD-4 to KFD-5 to
KFD-6 transition.

We make five contributions:

\begin{enumerate}[leftmargin=*]
  \item We define an \episode{} historical asset and distinguish it from a
        transcript, event log, provenance graph, embedding index, and endpoint
        state.
  \item We specify a staged discovery mechanism that keeps perspective replay,
        candidate genesis, qualification, held-out evaluation, and promotion
        authority independently inspectable.
  \item We explain why control of historical representation affects the power
        to define future questions, while deriving stewardship obligations
        rather than a command to retain everything.
  \item We freeze one implementation-backed KFD evidence cut and grade what it
        supports: three live cases and six tracks, with three bounded
        software-domain acceptances and three provisional candidates.
  \item We propose prospective ablation, shadow-loop, transfer, and governance
        evaluations capable of falsifying the stronger thesis.
\end{enumerate}

The paper does not claim autonomous Primitive discovery, universal ontology,
historical novelty of current KFD candidates, or a quantified return on
Episode retention. It claims that a mechanism can be specified, its early
feasibility can be separated from its missing gates, and the next experiment
can be made answerable to evidence.

\section{Problem Statement and Research Questions}
\label{sec:problem}

Let an operational system at time $n$ act with ontology $O_n$. Its capture
schema determines which entities, relations, state changes, and responsibilities
are directly representable. A fixed-ontology history $H(O_n)$ can be perfectly
internally consistent while remaining unable to express a distinction $x$
that becomes load-bearing later.

An ontology revision is justified only when the new distinction changes a
required decision and cannot be recovered from the old representation at lower
total burden. The problem is therefore not to maximize labels. It is to detect
when repeated reconstruction, hidden coordination, divergent consequences, or
causal variables indicate that $O_n$ is missing an independently addressable
responsibility.

We study four research questions:

\begin{description}[leftmargin=2.1em,style=nextline]
  \item[RQ1: Representation.] Which properties distinguish an Episode corpus
  that preserves future ontology-revision capacity from anonymous traces or
  endpoint facts?
  \item[RQ2: Mechanism.] Can declared perspective replay improve candidate
  genesis while KFD-5 qualification prevents narrative or method artifacts
  from becoming Primitives?
  \item[RQ3: Evaluation.] Under shared evidence and resource budgets, does an
  Episode-based loop outperform fixed-ontology and no-new-Primitive baselines
  on held-out reconstruction, transfer, and false-candidate rejection?
  \item[RQ4: Governance.] Which integrity, retention, deletion, consent,
  portability, access, and promotion boundaries are necessary when a corpus
  may influence not only answers but future problem definitions?
\end{description}

These questions deliberately separate technical possibility from legitimate
authority. A system may generate a useful candidate and still lack the right
to retain the underlying experience, represent a participant, or change an
adopted ontology.
